\documentclass[11pt]{article}
% Shared local declarations. Document roots use the repository's business-admin family.
\usepackage{business-admin}
\usepackage{amsmath,amssymb,tabularx,longtable,array,booktabs,enumitem}
\usepackage{ragged2e}
\newcolumntype{Y}{>{\RaggedRight\arraybackslash}X}
\newcommand{\SuiteMeta}[2]{\def\DocTitle{#1}\def\DocSubtitle{#2}\def\DocVersion{1.0}\def\DocStatus{Educational and operational reference}\def\DocOwner{Parliamentary Procedure and Governance Suite}\def\DocSummary{Independently authored reference and working materials, based on RONR 12th edition.}}
\newcommand{\SuiteOpen}{\makedoctitle\vspace{-0.5em}\noindent\rule{\linewidth}{0.4pt}\par\medskip}
\newcommand{\RonrAnchor}[1]{\par\smallskip{\footnotesize\textbf{Primary reading:} Robert's Rules of Order Newly Revised, 12th ed., #1. Original synthesis; not official interpretation.}\par}
\newcommand{\WriteLine}{\par\noindent\rule{\linewidth}{0.35pt}\par\vspace{0.5cm}}
\newcommand{\WriteLines}[1]{\par\noindent\rule{\linewidth}{0.35pt}\par\vspace{0.35cm}\noindent\rule{\linewidth}{0.35pt}\par\vspace{0.35cm}\noindent\rule{\linewidth}{0.35pt}\par\vspace{0.25cm}}
\newcommand{\Field}[1]{\par\medskip\noindent\textbf{#1}\par\WriteLines{3}}
\newcommand{\GovCheckbox}{\(\square\)}
\setlength{\emergencystretch}{3em}

\SuiteMeta{Operating Procedure: Handle a Point of Order or Appeal}{SOP | controlled meeting administration}
\begin{document}
\SuiteOpen
\RonrAnchor{23; 24; 62}
\section{Trigger and accountability}
\textbf{Trigger:} On alleged procedural breach\quad\textbf{Responsible:} Chair / appellant / secretary

\textbf{Primary reading:} RONR (12th ed.) 23; 24; 62

\section{Operating steps}
\begin{enumerate}[leftmargin=*,itemsep=4pt]
\item Stop only as necessary for timely procedural question.
\item Ask member to identify the violation.
\item Rule based on governing authority.
\item If proper appeal is made and seconded, state the appeal question.
\item Apply correct debate conditions and determine whether ruling stands.
\end{enumerate}

\section{Quality gate}
The secretary or designated independent reviewer must be able to trace the action to a valid convening authority, applicable governing rule, eligible voters where relevant, the chosen disposition and an official record. A process pass does not establish statutory legality.

\section{Exception and escalation}
The matter invokes continuing breach or chair abuse remedies. Document the issue, stop any unauthorized substantive action and obtain qualified parliamentary or legal guidance as appropriate.

\section{Retention and audit evidence}
\begin{itemize}[leftmargin=*,itemsep=4pt]
\item Approved meeting notice and governing authority where applicable.
\item Attendance and quorum evidence when a vote is involved.
\item Exact action text, vote basis and final disposition.
\item Corrections, chair rulings and approvals needed to reconstruct the transaction.
\end{itemize}
\end{document}
